\section{Multi Company}

\begin{frame}[fragile]{What Is Multi-Company?}
	\begin{itemize}
		\item Multi-company allows one database to serve multiple companies.
		\item Records often belong to a \texttt{company\_id}.
		\item Users may have access to one or several companies.
\begin{block}{Typical Field}
\begin{lstlisting}
company_id = fields.Many2one(
	'res.company',
	required=True,
	default=lambda self: self.env.company,
)
\end{lstlisting}
\end{block}
	\end{itemize}
\end{frame}

\begin{frame}[fragile]{Company Field on Business Models}
\begin{lstlisting}
class Student(models.Model):
	_name = 'student.student'

	name = fields.Char(required=True)
	company_id = fields.Many2one(
		'res.company',
		string='Company',
		required=True,
		index=True,
		default=lambda self: self.env.company,
	)
\end{lstlisting}
	\begin{itemize}
		\item Default to \texttt{env.company}.
		\item Index \texttt{company\_id} because security domains use it often.
	\end{itemize}
\end{frame}

\begin{frame}[fragile]{Multi-Company Record Rule}
\begin{lstlisting}[style=xml]
<record id="student_comp_rule" model="ir.rule">
  <field name="name">Student multi-company</field>
  <field name="model_id" ref="model_student_student"/>
  <field name="domain_force">[('company_id', 'in', company_ids)]</field>
</record>
\end{lstlisting}
	\begin{itemize}
		\item \texttt{company\_ids} = companies currently allowed in the user environment.
		\item This is the standard isolation pattern.
	\end{itemize}
\end{frame}

\begin{frame}[fragile]{Environment Helpers for Company}
\begin{lstlisting}
# Active company
company = self.env.company

# Allowed companies
companies = self.env.companies

# Force operations in another company
records = self.env['student.student'].with_company(other_company).search([])
\end{lstlisting}
	\begin{itemize}
		\item Prefer \texttt{with\_company()} over hard-coding company IDs.
	\end{itemize}
\end{frame}

\begin{frame}[fragile]{Views and Company Fields}
\begin{lstlisting}[style=xml]
<field name="company_id" groups="base.group_multi_company"/>
\end{lstlisting}
	\begin{itemize}
		\item Often show \texttt{company\_id} only when multi-company is enabled.
		\item In many2one fields, domain can restrict to current companies:
\begin{lstlisting}[style=xml]
<field name="classroom_id"
	   domain="[('company_id', 'in', company_ids)]"/>
\end{lstlisting}
	\end{itemize}
\end{frame}

\begin{frame}[fragile]{Check Company Consistency}
\begin{lstlisting}
@api.constrains('classroom_id', 'company_id')
def _check_company_consistency(self):
	for student in self:
		if student.classroom_id and student.classroom_id.company_id != student.company_id:
			raise ValidationError(
				"Student and Classroom must belong to the same company."
			)
\end{lstlisting}
	\begin{itemize}
		\item Related records should usually share the same company.
	\end{itemize}
\end{frame}

\begin{frame}{Multi-Company Best Practices}
	\begin{itemize}
		\item Add \texttt{company\_id} early on shared business models.
		\item Add a company record rule.
		\item Default from \texttt{env.company}.
		\item Validate consistency between related records.
		\item Test with users who can access more than one company.
	\end{itemize}
\end{frame}
